% Kandydat restrukturyzacji; nie jest aktywnym rozdziałem Atlasu.
\section{Transformacja Fouriera funkcji radialnych}
\providecommand{\atlasref}[2]{#2}

\subsection{Transformacja Fouriera funkcji radialnych}

Niech $f(x)=F(|x|)$ będzie funkcją radialną w $L^1(\R^d)$.
Niezmienniczość względem ortogonalnych zmian współrzędnych
w Stwierdzeniu~\ref{prop:dt-fourier-l1-rules} daje
$\widehat f(O\xi)=\widehat f(\xi)$ dla $O\in O(d)$.
Można zatem zastąpić całkę wielowymiarową całką po
promieniu. W wymiarze jeden otrzymujemy bezpośrednio
\[
 \widehat f(\xi)=\sqrt{2/\pi}\int_0^\infty
 F(r)\cos(r\xi)\,\mathrm d\lambda_1(r).
\]

\begin{lemat}[Całka po sferze]
\label{lem:dt-fourier-sphere-kernel}
Niech $d\ge2$ i $\sigma_{d-1}$ będzie miarą powierzchniową
na sferze jednostkowej $S^{d-1}$, czyli miarą Hausdorffa
o wymiarze $d-1$. Dla $t\in\R$ mamy
\[
 \int_{S^{d-1}}e^{-it\omega_1}\,\mathrm d\sigma_{d-1}(\omega)
 =|S^{d-2}|\int_{-1}^1 e^{-itu}(1-u^2)^{(d-3)/2}
 \,\mathrm d\lambda_1(u).
\]
Dla $d=2$ przyjmujemy $|S^0|=2$.
\end{lemat}
\begin{proof}
Poza dwoma biegunami stosujemy parametryzację
$\omega=(u,\sqrt{1-u^2}\,\eta)$, $\eta\in S^{d-2}$,
$-1<u<1$. Wektor pochodnej względem $u$ ma normę
$(1-u^2)^{-1/2}$ i jest prostopadły do kierunków
stycznych w $\eta$, ponieważ te ostatnie są prostopadłe
do $\eta$. Długości w $d-2$ kierunkach stycznych
skaluje czynnik $\sqrt{1-u^2}$. Z wzoru na pole
powierzchni iloczyn czynników wynosi
$(1-u^2)^{(d-3)/2}$. Bieguny mają miarę zero.
Dla $d=2$ parametryzacja składa się z dwóch łuków
odpowiadających dwóm punktom $S^0$, a ten sam
rachunek daje element długości $(1-u^2)^{-1/2}$.
Całkowana funkcja zależy wyłącznie od $u$, co daje tezę.
\end{proof}

\begin{definicja}[Funkcja Bessela pierwszego rodzaju]
\label{def:dt-bessel-j}
Dla $\alpha\ge0$ i $t>0$ definiujemy
\[
 J_\alpha(t)=\sum_{m=0}^\infty
 \frac{(-1)^m(t/2)^{2m+\alpha}}
      {m!\,\Gamma(m+\alpha+1)}.
\]
Tutaj $\Gamma(a)=\int_0^\infty
s^{a-1}e^{-s}\,\mathrm d\lambda_1(s)$ dla $a>0$.
Stosunek modułów kolejnych składników szeregu wynosi
$t^2/(4(m+1)(m+\alpha+1))$, ponieważ całkowanie
przez części daje $\Gamma(a+1)=a\Gamma(a)$.
Szereg zatem zbiega bezwzględnie i jednostajnie na
zwartych przedziałach dodatnich $t$.
\end{definicja}

\begin{twierdzenie}[Wzór radialny z funkcją Bessela]
\label{thm:dt-fourier-radial}
Dla $d\ge2$, $\alpha=d/2-1$ oraz $f(x)=F(|x|)\in L^1$,
przy $\rho=|\xi|>0$ zachodzi
\[
 \widehat f(\xi)=\rho^{-\alpha}\int_0^\infty
 F(r)J_\alpha(r\rho)r^{\alpha+1}\,\mathrm d\lambda_1(r).
\]
Przy $\rho=0$ wartość wynosi
$c_d|S^{d-1}|\int_0^\infty F(r)r^{d-1}\,\mathrm d\lambda_1(r)$.
\end{twierdzenie}
\begin{proof}
Najpierw obliczymy jądro sferyczne. Dla $a,b>0$
podstawienie $s=vw$, $t=v(1-w)$ w iloczynie całek
definiujących $\Gamma(a)\Gamma(b)$ ma jakobian $v$.
Twierdzenie Tonellego daje wzór beta:
\[
 \int_0^1 w^{a-1}(1-w)^{b-1}\,\mathrm d\lambda_1(w)
 =\frac{\Gamma(a)\Gamma(b)}{\Gamma(a+b)}.
\]
Ponadto całka Gaussa daje $\Gamma(1/2)=\sqrt{\pi}$,
a całkowanie $e^{-|x|^2}$ we współrzędnych biegunowych
w $\R^n$ daje
$\pi^{n/2}=|S^{n-1}|\Gamma(n/2)/2$.
W szczególności $|S^{d-2}|=2\pi^{(d-1)/2}/\Gamma((d-1)/2)$.
Rozwijamy wykładnik w Lemacie~\ref{lem:dt-fourier-sphere-kernel}
w szereg, który jest jednostajnie bezwzględnie zbieżny
dla $-1\le u\le1$ przy ustalonym $t$.
Składniki nieparzyste całkują się do zera, a podstawienie
$w=u^2$ i wzór beta dają
\[
 \int_{-1}^1 u^{2m}(1-u^2)^{(d-3)/2}\,\mathrm d\lambda_1(u)
 =\frac{\Gamma(m+1/2)\Gamma((d-1)/2)}{\Gamma(m+d/2)}.
\]
Z rekurencji funkcji gamma
$\Gamma(m+1/2)=\sqrt{\pi}(2m)!/(4^m m!)$.
Podstawiając do szeregu i mnożąc przez $|S^{d-2}|$,
otrzymujemy dla $t>0$
\[
 \int_{S^{d-1}}e^{-it\omega_1}\,\mathrm d\sigma_{d-1}
 =(2\pi)^{d/2}t^{-\alpha}J_\alpha(t).
\]
Istotnie, po podstawieniu definicji $J_\alpha$ po obu
stronach współczynnik przy $(-1)^mt^{2m}$ jest równy
$2\pi^{d/2}/(4^m m!\Gamma(m+d/2))$.

Wracamy do transformacji $f$. Obracamy $\xi$ na
$\rho e_1$ i stosujemy współrzędne biegunowe:
\[
 \widehat f(\xi)=c_d\int_0^\infty F(r)r^{d-1}
 \left(\int_{S^{d-1}}e^{-ir\rho\omega_1}
 \,\mathrm d\sigma_{d-1}(\omega)\right)\mathrm d\lambda_1(r).
\]
Całkowalność wynika z $f\in L^1$ i modułu jądra równego
jedności. Po wstawieniu obliczonej całki sferycznej
czynniki $(2\pi)^{d/2}$ i $c_d$ skracają się, a
$d-1-\alpha=\alpha+1$. Daje to wzór.
Wartość w zerze wynika wprost z całki definiującej
transformatę; nie trzeba dzielić przez $\rho=0$.
\end{proof}

\begin{przyklad}[Transformata kuli]
\label{ex:dt-fourier-ball}
Dla $R>0$ i $d\ge2$ niech $f=\chi_{\{|x|\le R\}}$.
Z Twierdzenia~\ref{thm:dt-fourier-radial} wynika dla
$\rho=|\xi|>0$
\[
 \widehat f(\xi)=R^{d/2}\rho^{-d/2}J_{d/2}(R\rho).
\]
Aby obliczyć całkę radialną, różniczkujemy szereg
definiujący $t^{\alpha+1}J_{\alpha+1}(t)$ na zwartych
przedziałach dodatnich $t$. Zbieżność jednostajna
szeregu pochodnych wynika z kryterium ilorazowego.
Korzystając z $\Gamma(m+\alpha+2)=(m+\alpha+1)\Gamma(m+\alpha+1)$,
dostajemy
$(t^{\alpha+1}J_{\alpha+1}(t))'=t^{\alpha+1}J_\alpha(t)$.
Podstawienie $t=r\rho$ i granica w zerze (równa zero
z szeregu) dają wskazany wzór. W szczególności dla $d=2$
jest to $R J_1(R\rho)/\rho$. W zerze transformata jest
równa $c_d\lambda_d(B(0,R))$, zgodnie z definicją.
\end{przyklad}
